\section{Доказателство за баланса на потоците в ММО (Фиг. 9.8)}
\begin{figure}
    \centering
	\begin{tikzpicture}[
		>=Stealth, 
		node distance=2cm and 2.5cm,
		queue/.style={circle, draw, minimum size=1cm, thick},
		label_node/.style={font=\small}
		]
		
		% Nodes
		\node[queue] (n1) {1};
		\node[queue, above right=0.5cm and 2cm of n1] (n2) {2};
		\node[queue, below right=0.5cm and 2cm of n1] (n3) {3};
		\node[queue, right=4.5cm of n1] (n4) {4};
		
		% External Inflow arrows
		\draw[<-] (n1.west) -- ++(-1,0) node[left] {6};
		\draw[<-] (n2.north) -- ++(0,0.8) node[above] {2};
		\draw[<-] (n3.south) -- ++(0,-0.8) node[below] {4};
		
		% Outflow arrow
		\draw[->] (n4.east) -- ++(1,0);
		
		% Internal flow arrows with probabilities
		\draw[->] (n1.east) -- (n2.west) node[midway, above left] {0.7};
		\draw[->] (n1.east) -- (n3.west) node[midway, below left] {0.3};
		\draw[->] (n2.east) -- (n4.west);
		\draw[->] (n3.east) -- (n4.west);
		
		% Feedback Loop
		\draw[->] (n4.north) -- ++(0,3) -| (n1.north) 
		node[pos=0.25, above] {0.4};
				
		% Subscripts (Lambdas)
		\node[below=0.1cm of n1] {$\lambda_1$};
		\node[above right=0.1cm of n2] {$\lambda_2$};
		\node[below right=0.1cm of n3] {$\lambda_3$};
		\node[above right=0.1cm of n4] {$\lambda_4$};
				
	\end{tikzpicture}
    \caption{Фиг. 9.8}
\end{figure}
За изследване на мрежата на Джексън използваме закона за запазване на потоците, съгласно който сумарната интензивност на входящите в даден възел потоци е равна на сумарната интензивност на изходящите потоци.

\subsection{Параметри на системата}
Векторът на външните входящи потоци е $a = [6, 2, 4, 0]$. Матрицата на вероятностите за преход $P$ е:
\[
P = \begin{pmatrix} 
0 & 0.7 & 0.3 & 0 \\ 
0 & 0 & 0 & 1 \\ 
0 & 0 & 0 & 1 \\ 
0.4 & 0 & 0 & 0 
\end{pmatrix}
\]

\subsection{Уравнения на трафика}
Съгласно общия вид $\lambda_{j} = a_{j} + \sum_{i=1}^{k} \lambda_{i} p_{ij}$, съставяме системата:
\begin{itemize}
    \item За възел 1: $\lambda_1 = 6 + 0.4\lambda_4$ 
    \item За възел 2: $\lambda_2 = 2 + 0.7\lambda_1$ 
    \item За възел 3: $\lambda_3 = 4 + 0.3\lambda_1$ 
    \item За възел 4: $\lambda_4 = \lambda_2 + \lambda_3$ 
\end{itemize}

\subsection{Решение на системата}
Чрез заместване на $\lambda_2$ и $\lambda_3$ в уравнението за $\lambda_4$:
\[ \lambda_4 = (2 + 0.7\lambda_1) + (4 + 0.3\lambda_1) = 6 + \lambda_1 \]
Замествайки в уравнението за $\lambda_1$:
\[ \lambda_1 = 6 + 0.4(6 + \lambda_1) \Rightarrow 0.6\lambda_1 = 8.4 \Rightarrow \lambda_1 = 14 \]
Следователно получените интензивности са:
\[ \lambda = [14, 11.8, 8.2, 20]^T \]

\subsection{Проверка на баланса}
Общият входящ поток в мрежата е:
\[ a_{total} = \sum a_i = 6 + 2 + 4 + 0 = 12 \]
Общият изходящ поток напуска системата през възел 4. Тъй като вероятността за обратна връзка от възел 4 към възел 1 е $p_{41} = 0.4$, вероятността за напускане е $1 - 0.4 = 0.6$.
\[ \lambda_{out} = \lambda_4 \times (1 - p_{41}) = 20 \times 0.6 = 12 \]
Тъй като $a_{total} = \lambda_{out}$, законът за запазване на потоците е изпълнен.

\section{Изчисляване на вероятностите за състоянието на ММО (Фиг. 9.10)}
\begin{figure}
    \centering
    \begin{tikzpicture}[
        >=Stealth, 
        node distance=3cm,
        state/.style={circle, draw, minimum size=1cm, thick},
        label_node/.style={font=\small}
        ]
        
        % Nodes
        \node[state] (s1) {1};
        \node[state, right of=s1] (s2) {2};
        \node[state, right of=s2] (s3) {3};
        
        % External Inflow
        \draw[<-] (s1.west) -- ++(-0.8,0) node[above left] {22};
        
        % Transitions between 1 and 2
        \draw[->] (s1.north east) [bend left=45] to node[above] {1} (s2.north west);
        \draw[->] (s2.south west) [bend left=45] to node[below] {0.2} (s1.south east);
        
        % Transitions between 2 and 3
        \draw[->] (s2.north east) [bend left=45] to node[above] {0.8} (s3.north west);
        \draw[->] (s3.south west) [bend left=45] to node[below] {0.5} (s2.south east);
        
        % External Outflow
        \draw[->] (s3.east) -- ++(0.8,0) node[above left] {0.5};
        
    \end{tikzpicture}
    \caption{Фиг. 9.10}
\end{figure}
За да определим вероятността във възлите 1 и 3 да има съответно 2 и 3 заявки, използваме принципа на декомпозиция на Джексън. Съгласно него, състоянието на мрежата се разглежда като произведение от състоянията на отделните СМО.

\subsection{Изходни параметри и интензивности}
Въз основа на анализа на трафика за Фиг. 9.10 са определени следните стойности:
\begin{itemize}
    \item Интензивност на входящия поток във възел 1: $\lambda_1 = 33$ заявки/час.
    \item Интензивност на обслужване във възел 1 (едноканален): $\mu_1 = 40$ заявки/час.
    \item Интензивност на входящия поток във възел 3: $\lambda_3 = 44$ заявки/час.
    \item Интензивност на обслужване във възел 3 (едноканален): $\mu_3 = 50$ заявки/час.
\end{itemize}

\subsection{Коефициенти на натоварване ($\rho$)}
Коефициентите на натоварване за съответните фази се пресмятат по формулата $\rho_i = \frac{\lambda_i}{\mu_i}$:
\begin{itemize}
    \item За възел 1: $\rho_1 = \frac{33}{40} = 0.825$.
    \item За възел 3: $\rho_3 = \frac{44}{50} = 0.88$.
\end{itemize}

\subsection{Пресмятане на индивидуалните вероятности}
Тъй като и двата възела са едноканални, използваме формулата за вероятност за $s$ заявки: $p_i(s) = \rho_i^s (1 - \rho_i)$:

\begin{itemize}
    \item \textbf{Вероятност във възел 1 да има 2 заявки:}
    \[ p_1(2) = 0.825^2 \times (1 - 0.825) = 0.680625 \times 0.175 \approx 0.119 \text{ } \]
    
    \item \textbf{Вероятност във възел 3 да има 3 заявки:}
    \[ p_3(3) = 0.88^3 \times (1 - 0.88) = 0.681472 \times 0.12 \approx 0.0818 \text{ } \]
\end{itemize}

\subsection{Обща вероятност за системата}
Общата вероятност за едновременното наличие на посочения брой заявки във възли 1 и 3 е:
\[ p(s_1=2, s_3=3) = p_1(2) \times p_3(3) \]
\[ p(2, 3) = 0.119 \times 0.0818 \approx 0.00973 \]